\chapter{METODE DAN RANCANGAN PENELITIAN}

\section{Deskripsi Umum Penelitian}

Penelitian ini menggunakan pendekatan eksperimental untuk mengukur hubungan antara kualitas ringkasan dan efisiensi inferensi. Model yang diuji adalah Qwen3-0.6B dengan sebuah router ringan yang memilih kelompok lapisan untuk dijalankan berdasarkan prompt. Tugasnya adalah meringkas berita berbahasa Indonesia pada subset Indonesia dari XL-Sum. Hasil router akan dibandingkan dengan model penuh dan pelompatan lapisan statis.

Pelatihan dilakukan dalam dua tahap. Pada tahap pertama, seluruh parameter Qwen dibekukan dan hanya parameter router yang dilatih. Pada tahap kedua, router dibekukan dan adapter LoRA dilatih untuk mengompensasi perubahan kualitas akibat pelompatan lapisan, mengikuti pola dua tahap FiRST. Router memilih satu pola kelompok lapisan untuk setiap prompt; pola itu digunakan selama sisa pemrosesan prompt dan semua langkah pembentukan token keluaran. Dampaknya dinilai dari kualitas ringkasan dan waktu inferensi aktual, bukan hanya dari jumlah lapisan yang dilewati.

\section{Model dan Data Penelitian}

\subsection{Model Qwen3-0.6B}

Model utama yang digunakan adalah \texttt{Qwen/Qwen3-0.6B}. Dokumentasi model mencatat bahwa Qwen3-0.6B merupakan model bahasa kausal dengan 28 lapisan \citep{qwen2025modelcard}. Lapisan dibagi menjadi tujuh kelompok berurutan, masing-masing berisi empat lapisan. Jika lapisan diberi nomor dari satu, kelompok $G_i$ mencakup lapisan $4i-3$ hingga $4i$, untuk $i=1,\ldots,7$. Kelompok pertama ($G_1$) dan terakhir ($G_7$) selalu dijalankan. Router hanya menentukan apakah lima kelompok tengah ($G_2$ sampai $G_6$) dijalankan atau dilewati. Menjalankan $G_1$ dan $G_7$ merupakan asumsi rancangan; penelitian ini tidak menganggap keduanya telah terbukti sebagai kelompok terpenting.

Revisi bobot model, tokenizer, dan versi perangkat lunak akan dicatat agar konfigurasi eksperimen dapat diketahui. Generasi dilakukan dalam mode non-thinking, yang mencegah Qwen menghasilkan teks penalaran sebelum jawaban akhir. Jendela konteks (\textit{context window}), yaitu kapasitas token gabungan prompt dan keluaran dalam satu urutan, dibatasi hingga 4.096 token. Batas ini adalah pengaturan eksperimen, bukan kapasitas maksimum bawaan model. Paling banyak 256 token digunakan untuk keluaran melalui parameter \texttt{max\_new\_tokens=256}; karena itu, prompt dibatasi hingga 3.840 token, termasuk instruksi dan token format chat. Panjang prompt dihitung dengan tokenizer Qwen3-0.6B setelah format chat diterapkan. Model dapat berhenti sebelum menghasilkan 256 token. Batas prompt dan keluaran yang sama berlaku bagi semua metode.

Generasi menggunakan pemilihan token greedy (\textit{greedy decoding}): pada setiap langkah, model memilih token dengan skor tertinggi tanpa mengambil sampel secara acak. Pengaturan yang digunakan adalah \texttt{do\_sample=False} dan \texttt{num\_beams=1}. Generasi berhenti saat model menghasilkan token akhir respons sesuai konfigurasi model atau saat mencapai 256 token keluaran. Token akhir respons termasuk dalam batas tersebut; panjang minimum keluaran tidak ditetapkan. Aturan penghentian yang sama berlaku bagi semua metode. Identitas token akhir respons akan dicatat bersama revisi konfigurasi model. Tidak dilakukan variasi tambahan pada pengaturan generasi antarmetode.

\subsection{XL-Sum Subset Indonesia}

Data akan diambil dari dataset XL-Sum di Hugging Face dengan konfigurasi \texttt{indonesian} \citep{xlsumDatasetCard}. Repositori resmi XL-Sum mencantumkan 38.242 pasangan untuk data latih, 4.780 untuk data validasi, dan 4.780 untuk data uji \citep{xlsumRepository}. Setelah pembersihan, penyaringan, dan penghapusan duplikasi, penelitian menggunakan 1.000 artikel latih, 200 artikel validasi, dan 300 artikel uji yang dipilih secara acak dengan seed 42.

Artikel menjadi masukan untuk prompt peringkasan. Ringkasan acuan digunakan untuk menilai kualitas keluaran dan untuk membentuk target pelatihan router pada data latih. Data dipisahkan menjadi data latih (\textit{train}), validasi (\textit{validation}), dan uji (\textit{test}). Data uji tidak akan digunakan untuk menentukan target pelatihan ataupun memilih konfigurasi.

\medskip
\noindent\textbf{Pembersihan Teks dan Duplikasi}\par\smallskip

Spasi dan baris kosong hanya dihapus dari awal dan akhir artikel serta ringkasan acuan. Huruf besar, tanda baca, spasi di antara kata, dan isi teks lainnya dipertahankan. Pasangan artikel--ringkasan dikeluarkan jika salah satu teks kosong setelah pembersihan. Aturan ini diterapkan pada semua split sebelum prompt disusun dan panjang token diperiksa. Jumlah pasangan kosong yang dikeluarkan akan dicatat. Pengolahan teks untuk metrik ROUGE akan ditentukan secara terpisah.

Artikel dianggap duplikat jika teksnya sama persis setelah pembersihan. Dalam satu split, kemunculan pertama menurut urutan asli dataset dipertahankan bersama ringkasan acuannya. Jika artikel yang sama muncul pada beberapa split, pasangan yang dipertahankan berasal dari data uji jika tersedia; jika tidak, data validasi diprioritaskan daripada data latih. Salinan pada split lain dikeluarkan tanpa dipindahkan. Pemeriksaan dilakukan sebelum pemilihan subset. Identitas dan jumlah pasangan yang dikeluarkan karena duplikasi dalam satu split maupun antarsplit akan dicatat.

\medskip
\noindent\textbf{Batas Panjang dan Pemilihan Sampel}\par\smallskip

Artikel dikeluarkan jika prompt lengkap yang berisi artikel utuh melebihi 3.840 token menurut tokenizer Qwen3-0.6B. Penghitungan mencakup instruksi dan format chat. Artikel tidak dipotong. Ringkasan acuan dikeluarkan jika panjangnya melebihi 256 token, termasuk token akhir respons. Ringkasan juga tidak dipotong. Penyaringan ini dilakukan pada data latih, validasi, dan uji; semua metode akan dievaluasi pada himpunan artikel uji yang sama. Jumlah pasangan yang dikeluarkan karena batas artikel, jumlah yang dikeluarkan karena batas ringkasan, dan jumlah akhir setiap split akan dicatat menggunakan format prompt final.

Gabungan prompt dan ringkasan acuan yang diberikan sebagai contoh untuk memandu prediksi model selama pelatihan (\textit{teacher forcing}) juga akan diperiksa agar tidak melebihi 4.096 token. Sampel dipilih secara acak tanpa pengembalian dari setiap split setelah pembersihan, penyaringan panjang, dan penghapusan duplikasi. Sebanyak 1.000 artikel digunakan untuk data latih, 200 artikel untuk data validasi, dan 300 artikel untuk data uji. Pemilihan menggunakan seed, yaitu angka awal untuk mengatur urutan bilangan acak, dengan nilai tetap 42. Identitas sampel disimpan, dan semua metode menggunakan subset evaluasi yang sama.

\medskip
\noindent\textbf{Format Prompt}\par\smallskip

Instruksi dan artikel dimasukkan sebagai pesan pengguna dengan format chat Qwen dalam mode non-thinking. Isi pesan yang digunakan adalah berikut; teks \texttt{\{artikel\}} diganti dengan artikel utuh.

\begin{quote}
Ringkas artikel berita berikut dalam bahasa Indonesia. Pertahankan informasi utama dan jangan menambahkan informasi yang tidak terdapat dalam artikel.

Artikel:\\
\texttt{\{artikel\}}
\end{quote}

Format yang sama digunakan saat menyiapkan data pelatihan router dan saat menghasilkan ringkasan dengan seluruh metode pembanding. Batas masukan dihitung setelah format chat diterapkan, termasuk penanda awal respons asisten. Artikel dengan prompt lengkap lebih dari 3.840 token dikeluarkan tanpa dipotong.

\section{Tahapan Penelitian}

Penelitian direncanakan melalui tahapan berikut.

\begin{enumerate}
  \item \textbf{Perumusan rancangan.} Kajian literatur, batasan, dan protokol evaluasi akan dilengkapi. Luaran tahap ini adalah spesifikasi eksperimen yang membedakan keputusan yang telah ditetapkan dari parameter yang masih harus dipilih menggunakan data validasi.
  \item \textbf{Persiapan data dan model.} Backbone dan split data akan disiapkan, prompt akan dibentuk, serta revisi model dan konfigurasi dicatat. Luaran tahap ini adalah proses masukan--keluaran model penuh yang dapat dijalankan ulang.
  \item \textbf{Implementasi eksekusi kondisional.} Router, mekanisme pelompatan kelompok, dan metode pembanding akan diimplementasikan. Kebenaran eksekusi, konsistensi jalur, dan penanganan KV cache akan diperiksa.
  \item \textbf{Pelatihan router dan LoRA.} Target keputusan router akan dibuat dari data latih melalui pencarian bertahap berbasis kerugian prediksi token dan proporsi kelompok aktif. Router dilatih dengan parameter Qwen dibekukan. Setelah konfigurasi router dipilih menggunakan data validasi, bobot router dibekukan dan adapter LoRA dilatih untuk model penuh, static skipping, dan routing adaptif. Tahap LoRA untuk routing adaptif turut menjaga tingkat skipping melalui penalti pada skor aktif router.
  \item \textbf{Evaluasi dan analisis.} Konfigurasi terpilih akan diuji pada data uji. Kualitas ringkasan, waktu inferensi, dan pola jalur yang dipilih akan dianalisis dan dibandingkan. Luaran tahap ini adalah hasil perbandingan beserta keterbatasannya.
  \item \textbf{Penulisan laporan.} Rancangan, implementasi, dan temuan akan didokumentasikan. Penulisan dilakukan sejalan dengan tahapan penelitian.
\end{enumerate}

Pemeriksaan awal implementasi merupakan bagian dari eksperimen utama, bukan tahap pilot yang terpisah. Saat draf ini disusun, topik, model, dataset, dan metode pembanding telah ditetapkan. Skrip inferensi model penuh sudah tersedia; implementasi routing dan evaluasi penelitian masih akan dikerjakan.

\section{Perancangan Sistem}

\subsection{Arsitektur Umum}

Gambar~\ref{fig:arsitektur-routing} menunjukkan alur sistem. Prompt diproses terlebih dahulu melalui kelompok pertama. Dari keluaran kelompok ini, router menerima rata-rata vektor representasi internal (\textit{hidden state}) pada posisi token artikel. Posisi token instruksi, token khusus, dan padding tidak disertakan dalam rata-rata. Router kemudian menghasilkan lima keputusan biner, satu untuk setiap kelompok tengah. Kelompok yang dijalankan tetap mengikuti urutan lapisan pada Qwen.

\begin{figure}[H]
  \centering
  \resizebox{\textwidth}{!}{% Diagram arsitektur sequence-level routing.
\begin{tikzpicture}[
  node distance=0.68cm and 1.2cm,
  block/.style={draw, rounded corners, align=center, text width=6.3cm,
    minimum height=0.78cm, inner sep=6pt, font=\small},
  fixed/.style={block, fill=gray!10},
  conditional/.style={block, fill=blue!6},
  router/.style={block, text width=5.1cm, fill=orange!10},
  flow/.style={-{Latex}, thick},
  control/.style={-{Latex}, dashed, thick}
]
  % Main model path
  \node[block] (input) {Artikel dan instruksi peringkasan};
  \node[fixed, below=of input] (g1) {$G_1$: lapisan 1--4\\Selalu dijalankan};
  \node[conditional, below=of g1] (middle) {Kelompok tengah $G_2$--$G_6$\\Tiap kelompok dijalankan ($m_i=1$) atau dilewati ($m_i=0$)};
  \node[fixed, below=of middle] (g7) {$G_7$: lapisan 25--28\\Selalu dijalankan};
  \node[fixed, below=of g7] (head) {Normalisasi akhir dan kepala bahasa (LM head)};
  \node[block, below=of head] (output) {Ringkasan keluaran};

  \foreach \a/\b in {input/g1,g1/middle,middle/g7,g7/head,head/output}
    \draw[flow] (\a.south) -- (\b.north);

  % Sequence-level router branch
  \node[router, right=of g1] (pool) {Rata-rata vektor representasi internal\\token artikel setelah $G_1$};
  \node[router, below=of pool] (mlp) {MLP $1024\mathbin{\to}64\mathbin{\to}5$\\ReLU dan 5 keluaran sigmoid};
  \node[router, below=of mlp] (mask) {Ambang mengubah keluaran sigmoid\\menjadi keputusan biner $G_2$--$G_6$};

  \draw[flow] (g1.east) -- (pool.west);
  \draw[flow] (pool.south) -- (mlp.north);
  \draw[flow] (mlp.south) -- (mask.north);
  \draw[control] (mask.west) -- ++(-0.6,0) |- (middle.east);

  \node[align=left, text width=5.1cm, font=\footnotesize, below=0.8cm of mask] (note)
    {Garis putus-putus menunjukkan keputusan router.\\Satu jalur dipilih per prompt untuk sisa prompt dan seluruh pembentukan keluaran.\\Tahap 1: Qwen dibekukan dan router dilatih.\\Tahap 2: router dibekukan dan LoRA dilatih pada attention serta FFN.};
\end{tikzpicture}
}
  \caption{Arsitektur konseptual sequence-level routing pada Qwen3-0.6B.}
  \label{fig:arsitektur-routing}
\end{figure}

Hasil kelompok pertama yang digunakan untuk menghitung representasi router juga dipakai untuk melanjutkan pemrosesan prompt melalui kelompok yang dipilih. Dengan demikian, kelompok pertama tidak dihitung ulang hanya untuk menjalankan router. Jalur yang dipilih dipertahankan untuk sisa pemrosesan prompt dan semua langkah pembentukan keluaran. Cache key dan value (KV cache), yaitu penyimpanan representasi yang digunakan kembali pada langkah inferensi berikutnya, disiapkan untuk lapisan yang aktif. Pada setiap token keluaran baru, kelompok pertama tetap dijalankan, tetapi router tidak memilih ulang jalur. Waktu untuk memperoleh representasi dan menjalankan router akan dimasukkan dalam waktu inferensi total. Integrasi cache dan kesetaraan perilaku dengan model penuh akan diperiksa dalam implementasi.

\subsection{Pemilihan Jalur serta Pelatihan Router dan LoRA}

Jalur adalah pola kelompok lapisan yang dijalankan untuk satu prompt. Router membuat keputusan berdasarkan prompt satu kali saja; keputusan itu tidak berubah untuk setiap token keluaran. Penelitian ini tidak menggunakan mekanisme penghentian lebih awal yang memutuskan pada setiap lapisan berdasarkan tingkat keyakinan model terhadap keluaran.

\medskip
\noindent\textbf{Kelompok Lapisan dan Keputusan Biner}\par\smallskip

Kelompok $G_1$ dan $G_7$ selalu dijalankan. Untuk tiap kelompok $G_i$, dengan $i=2,\ldots,6$, router menentukan nilai $m_i\in\{0,1\}$. Nilai satu berarti kelompok dijalankan dan nilai nol berarti kelompok dilewati. Pola lengkap ditulis sebagai $(1,m_2,m_3,m_4,m_5,m_6,1)$. Jumlah lapisan aktif tidak dipatok dan dapat berubah antarartikel dari 8 hingga 28, dengan perubahan sebanyak empat lapisan. Urutan lapisan yang dijalankan selalu mengikuti urutan Qwen.

\medskip
\noindent\textbf{Arsitektur Router dan Rencana Pelatihan}\par\smallskip

Router menggunakan jaringan perceptron multilapis (\textit{multilayer perceptron}, MLP), yaitu jaringan saraf yang memproses masukan melalui lapisan-lapisan terhubung. Masukan MLP adalah rata-rata representasi token artikel setelah $G_1$. MLP memiliki satu lapisan tersembunyi berisi 64 unit dengan fungsi aktivasi ReLU dan lima keluaran sigmoid. Setiap keluaran berpasangan dengan satu kelompok tengah. Nilai keluaran yang lebih besar atau sama dengan 0,5 diterjemahkan sebagai keputusan aktif, sedangkan nilai di bawah 0,5 diterjemahkan sebagai keputusan lewati. Saat inferensi, router hanya menerima representasi artikel dari prompt, bukan ringkasan acuannya.

Lima keputusan biner memungkinkan 32 kombinasi secara matematis. Namun, router tidak mengklasifikasikan 32 nomor jalur, dan pelatihan tidak harus menilai semua kombinasi tersebut. Pada tahap router, parameter Qwen tetap dibekukan, termasuk parameter kelompok pertama yang menghasilkan representasi router. Pada tahap berikutnya, router tetap dibekukan dan adapter LoRA dilatih pada modul attention dan FFN seluruh lapisan Qwen. Konfigurasi awal FiRST menggunakan rank 8, faktor skala 32, dan dropout 0,1; konfigurasi ini menjadi titik awal yang akan diverifikasi pada data validasi, bukan klaim konfigurasi optimal untuk Qwen3-0.6B.

Target keputusan dibuat secara luring (\textit{offline}) dari data latih. Untuk sebuah artikel dengan ringkasan acuan, setiap kandidat jalur dinilai menggunakan kerugian prediksi token dengan \textit{teacher forcing}, yaitu proses yang memberikan token acuan sebelumnya ketika menghitung prediksi token berikutnya. Skor kandidat jalur dirumuskan sebagai
\begin{equation}
  J(m) = \operatorname{NLL}(m) + \lambda C(m),
  \qquad C(m) = \frac{m_2+m_3+m_4+m_5+m_6}{5},
  \label{eq:route-objective}
\end{equation}
\noindent dengan $\operatorname{NLL}(m)$ sebagai rata-rata kerugian prediksi token pada jalur $m$, $C(m)$ sebagai proporsi kelompok tengah yang aktif, dan $\lambda$ sebagai bobot penghematan komputasi. Karena kelima kelompok memiliki jumlah lapisan yang sama, $C(m)$ digunakan sebagai ukuran sederhana proporsi kelompok yang dijalankan. Waktu inferensi aktual tetap diukur secara terpisah pada tahap evaluasi.

Pencarian target dimulai dari jalur penuh, yaitu semua kelompok tengah aktif. Pada setiap langkah, satu kelompok tengah yang masih aktif dicoba untuk dilewati. Kandidat dengan nilai $J(m)$ paling rendah dipilih jika nilainya lebih rendah daripada jalur saat ini. Proses berhenti ketika tidak ada penghapusan kelompok yang menurunkan nilai tersebut. Cara ini menghasilkan target biner untuk $G_2$ sampai $G_6$ tanpa kewajiban mengevaluasi seluruh 32 pola jalur.

Router dilatih menggunakan \textit{binary cross-entropy with logits} pada lima target biner tersebut. Penalti penghematan tidak ditambahkan ke fungsi klasifikasi karena penghematan telah digunakan saat membentuk target jalur. Konfigurasi awal router menggunakan optimizer AdamW dengan \textit{learning rate} $10^{-3}$, \textit{weight decay} $10^{-4}$, paling banyak 50 epoch, dan \textit{early stopping} dengan kesabaran lima epoch. Kandidat nilai penalti target ditetapkan sebagai $\lambda\in\{0, 0{,}01, 0{,}05, 0{,}10, 0{,}20\}$. Satu nilai $\lambda$ dipilih menggunakan data validasi dan digunakan pula saat memilih pola static. Seed pemilihan subset ditetapkan sebesar 42, sedangkan seed utama pelatihan router ditetapkan sebesar 43. Seed tambahan 44 dan 45 dapat digunakan untuk analisis ketahanan apabila sumber daya mencukupi; hasilnya akan dilaporkan terpisah dari hasil seed utama. Ringkasan acuan hanya digunakan untuk membentuk target pada pelatihan; saat inferensi, router hanya menggunakan representasi prompt. Data uji tidak digunakan untuk membentuk target atau memilih konfigurasi.

Setelah router terpilih, bobot router dibekukan dan LoRA dilatih menggunakan prompt serta ringkasan acuan dengan \textit{teacher forcing}. Adapter dipasang pada proyeksi attention dan FFN di seluruh lapisan, sedangkan bobot dasar Qwen tidak diperbarui. Konfigurasi awal FiRST menggunakan rank 8, faktor skala 32, dan dropout 0,1 \citep{jain2025first}; nilai ini menjadi titik awal, bukan klaim optimal untuk Qwen3-0.6B. Untuk metode routing adaptif, fungsi objektif LoRA menambahkan penalti non-skip berbobot $\beta$ pada probabilitas kelompok aktif, $L_{\mathrm{LoRA}}=L_{\mathrm{NLL}}+\beta\sum_{i=2}^{6}p_i$, dengan $p_i$ sebagai probabilitas router untuk menjalankan kelompok $G_i$. Penalti ini menjaga agar perubahan representasi akibat LoRA tidak menghilangkan perilaku skipping. Kandidat nilai $\beta$ serta rincian optimisasi LoRA yang belum ditetapkan dipilih atau diverifikasi menggunakan data validasi; data uji tidak digunakan untuk keputusan tersebut. Model penuh+LoRA dan static skipping+LoRA dilatih dengan konfigurasi adapter dan protokol yang sama, tanpa penalti non-skip untuk jalur yang ditetapkan tetap. Router tanpa LoRA menjadi kontrol ablation dari metode adaptif.

\subsection{Bagan Alir Metodologi}

Gambar~\ref{fig:bagan-alir-routing} merangkum urutan penelitian. Bagan membedakan persiapan yang telah tersedia dari pekerjaan implementasi dan pengujian yang masih direncanakan.

\begin{figure}[H]
  \centering
  \begin{tikzpicture}[
    node distance=0.35cm,
    block/.style={draw, rounded corners, align=center, text width=10.3cm, inner sep=6pt, font=\small},
    arrow/.style={-{Latex}, thick}
  ]
    \node[block] (scope) {Sudah: penetapan topik, model, dataset,\\pertanyaan penelitian, dan pembanding};
    \node[block, below=of scope] (basic) {Sudah: penyiapan skrip inferensi dasar\\Belum mengevaluasi metode routing};
    \node[block, below=of basic] (design) {Rencana: pelengkapan literatur, pengaturan router,\\pengolahan data, dan protokol evaluasi};
    \node[block, below=of design] (implementation) {Rencana: implementasi routing dan pembanding\\Luaran: eksekusi kondisional yang terverifikasi};
    \node[block, below=of implementation] (training) {Rencana: pelatihan router dan pemilihan konfigurasi\\Data latih untuk pelatihan; data validasi untuk pemilihan};
    \node[block, below=of training] (evaluation) {Rencana: evaluasi empat varian pada data uji\\Luaran: perbandingan kualitas, waktu, dan jalur};
    \node[block, below=of evaluation] (report) {Rencana: analisis, kesimpulan, dan laporan\\Penulisan berlangsung sepanjang penelitian};
    \draw[arrow] (scope) -- (basic);
    \draw[arrow] (basic) -- (design);
    \draw[arrow] (design) -- (implementation);
    \draw[arrow] (implementation) -- (training);
    \draw[arrow] (training) -- (evaluation);
    \draw[arrow] (evaluation) -- (report);
  \end{tikzpicture}
  \caption{Bagan alir metodologi dan status pekerjaan saat draf disusun.}
  \label{fig:bagan-alir-routing}
\end{figure}

\section{Rancangan Evaluasi}

\subsection{Metode Pembanding}

Evaluasi menggunakan empat varian: model penuh+LoRA, static skipping+LoRA, router+LoRA, dan router tanpa LoRA. Model penuh+LoRA menjadi acuan kualitas dan biaya tanpa pelompatan setelah adaptasi pada tugas. Pelompatan statis menggunakan satu pola tetap yang dipilih berdasarkan data validasi. Setelah dipilih, pola tersebut digunakan untuk semua artikel, selama sisa pemrosesan prompt, dan selama seluruh pembentukan keluaran. Router+LoRA menggunakan keputusan dari representasi prompt, sedangkan kontrol router-only menggunakan bobot backbone dan router dari tahap pertama tanpa kompensasi adapter. Data uji tidak digunakan untuk memilih pola statis atau konfigurasi router maupun LoRA.

Perbandingan kualitas utama akan mencakup router+LoRA terhadap full+LoRA dan static+LoRA. Perbandingan router+LoRA dengan router-only mengukur kontribusi tahap kompensasi LoRA. Setiap kandidat pola static dengan $G_1$ dan $G_7$ aktif dinilai pada data validasi menggunakan objective $J(m)$ pada Persamaan~\ref{eq:route-objective}. Pola dengan nilai objective paling rendah dipilih dan dikunci sebelum data uji digunakan. Jika terdapat nilai yang sama, pola dengan jumlah kelompok tengah aktif lebih sedikit dipilih. Jika masih sama, digunakan urutan leksikografis menaik pada $(m_2,m_3,m_4,m_5,m_6)$ dengan nilai 0 didahulukan daripada 1. Efisiensi utama dibandingkan menggunakan waktu inferensi aktual pada lingkungan dan pengaturan yang sama, tanpa penyetaraan biaya buatan berdasarkan jumlah lapisan. Jumlah lapisan aktif dan waktu inferensi aktual akan dicatat. Jumlah lapisan yang sama belum tentu menghasilkan biaya yang sama karena waktu routing, adapter, dan panjang keluaran dapat berbeda. Jumlah lapisan aktif tidak akan digunakan sebagai pengganti pengukuran jumlah operasi aritmetika bilangan titik mengambang (\textit{floating-point operations}, FLOPs) ataupun sebagai bukti percepatan. Waktu total mencakup pekerjaan router dan adapter yang benar-benar dijalankan; metode statis tidak dibebani biaya menjalankan router yang tidak digunakan.

\subsection{Pemeriksaan Kebenaran Implementasi}

Pemeriksaan awal akan memastikan bahwa ketika semua lapisan dijalankan, implementasi memberikan perilaku yang sesuai dengan model penuh pada kondisi yang sama. Pemeriksaan juga akan memastikan bahwa lapisan yang ditandai untuk dilewati memang tidak dieksekusi, jalur tidak berubah selama pembentukan keluaran, KV cache dikelola secara konsisten, bobot router tidak berubah pada tahap LoRA, dan adapter hanya dilatih pada tahap kedua. Pemeriksaan ini dilakukan sebelum hasil kualitas dan waktu ditafsirkan.

\subsection{Kualitas Ringkasan dan Efisiensi Inferensi}

Kualitas keluaran dibandingkan dengan ringkasan acuan menggunakan ROUGE-1, ROUGE-2, dan ROUGE-L. Metrik ini mengukur kecocokan kata tunggal (unigram), pasangan kata berurutan (bigram), dan subsekuens terpanjang yang sama. Penilaian dilakukan sepenuhnya secara otomatis; tidak ada penilaian manual atau pemeriksaan faktualitas manual pada sampel. F1 akan dilaporkan untuk ROUGE-1, ROUGE-2, ROUGE-L, dan BERTScore. Untuk setiap metrik, F1 dihitung per artikel lalu dirata-ratakan secara aritmetika dengan bobot yang sama bagi setiap artikel. Keempat rata-rata dilaporkan secara terpisah. F1 tidak dihitung dari presisi dan cakupan (\textit{recall}) yang terlebih dahulu dirata-ratakan.

ROUGE akan dihitung menggunakan implementasi \texttt{rouge-score} dengan tokenisasi bawaan, tanpa stemming. Tokenisasi tersebut menormalisasi huruf menjadi huruf kecil dan memisahkan tanda baca sesuai aturan implementasi, sedangkan teks sumber dan ringkasan yang disimpan tidak diubah. Metrik yang digunakan adalah \texttt{rouge1}, \texttt{rouge2}, dan \texttt{rougeL}; \texttt{rougeL} digunakan pada tingkat kalimat, bukan \texttt{rougeLsum}. Pengaturan ini dipilih agar tidak menerapkan stemmer bahasa Inggris pada teks bahasa Indonesia. Skor F1 dikumpulkan per artikel sebelum rata-rata makro dihitung.

BERTScore melengkapi ROUGE dengan membandingkan representasi kontekstual token pada ringkasan keluaran dan ringkasan acuan menggunakan kemiripan kosinus \citep{zhang2020bertscore}. Evaluasi akan menggunakan encoder \texttt{bert-base-multilingual-cased}, parameter bahasa \texttt{id}, dan lapisan ke-9 sesuai konfigurasi standar model tersebut pada pustaka BERTScore. IDF tidak digunakan (\texttt{idf=False}) sehingga semua token diberi bobot yang sama dan skor tidak bergantung pada korpus tambahan untuk menghitung frekuensi dokumen. Skor tidak disesuaikan dengan baseline (\texttt{rescale\_with\_baseline=False}) agar konfigurasi utama tidak bergantung pada berkas baseline tambahan yang dapat berbeda menurut versi perangkat lunak. Teks untuk BERTScore hanya menggunakan pembersihan awal--akhir yang telah ditetapkan; huruf besar, tanda baca, dan isi teks lainnya dipertahankan. Jika tokenisasi encoder menghasilkan lebih dari 510 token subkata, bagian akhir teks dipotong hingga batas tersebut sebelum penilaian, dan jumlah kasus pemotongan dicatat. Konfigurasi evaluator akan sama untuk semua metode. Metrik kecocokan terhadap acuan tidak menjamin faktualitas atau kualitas ringkasan secara keseluruhan. Forde et al. menemukan bahwa metrik otomatis dapat menjadi ukuran kualitas yang tidak andal pada peringkasan multibahasa, dengan masalah yang lebih besar pada bahasa Indonesia; karena itu, hasil di sini akan ditafsirkan terbatas pada metrik yang digunakan \citep{forde2024multilingeval}.

Selisih rata-rata F1 antar-metode akan disertai interval kepercayaan dua sisi 95\% menggunakan bootstrap berpasangan. Pada setiap pengambilan ulang (\textit{resampling}), artikel diambil dengan pengembalian sebanyak jumlah artikel uji semula. Indeks artikel yang sama digunakan untuk kedua metode dalam satu perbandingan, lalu selisih rata-rata F1 dihitung ulang. Penghitungan menggunakan skor per artikel yang telah disimpan, sehingga ringkasan tidak perlu dibuat ulang. Interval menggunakan metode persentil berdasarkan 10.000 pengambilan ulang dengan seed bootstrap 42; batas bawah dan atas adalah persentil 2,5\% dan 97,5\% dari distribusi selisih. Seed bootstrap ini terpisah dari seed pemilihan subset dan pelatihan router. Interval dihitung untuk router+LoRA dibandingkan dengan full+LoRA dan static+LoRA, serta router+LoRA dibandingkan dengan router-only sebagai ablation; semua perbandingan dihitung terpisah untuk keempat metrik. Hasil utama menggunakan satu seed pelatihan router, yaitu seed 43; variasi antarseed tidak menjadi bagian dari eksperimen utama. Interval yang dihitung dari artikel mengukur variasi antarartikel dan tidak mencakup variasi antarseed.

Waktu inferensi diukur per artikel dan mencakup waktu total, waktu pemrosesan prompt (\textit{prefill}), serta waktu pembentukan keluaran (\textit{decoding}). Batas awal dan akhir pengukuran setiap komponen akan ditetapkan agar biaya representasi serta router tercakup tanpa dihitung dua kali. Jumlah token keluaran juga dicatat: waktu yang lebih singkat dapat terjadi hanya karena model menghasilkan ringkasan yang lebih pendek. Sebelum pengukuran, setiap metode dijalankan tiga kali sebagai pemanasan dan hasilnya tidak dicatat. Setelah itu, pengukuran diulang lima kali pada setiap artikel. Median dari lima pengukuran digunakan sebagai waktu utama, sedangkan variasinya dilaporkan sebagai informasi tambahan.

Semua metode akan menggunakan perangkat, presisi numerik, prompt, dan pengaturan generasi yang sama. Mode non-thinking dan greedy decoding digunakan dengan \texttt{do\_sample=False}, \texttt{num\_beams=1}, dan batas keluaran 256 token. Hasil kualitas akan dianalisis bersama biaya inferensi aktual. Frekuensi setiap pola jalur yang dipilih akan dicatat untuk mengetahui apakah router hampir selalu memilih pola yang sama. Kondisi ini disebut \textit{routing collapse}. Temuan tersebut akan dilaporkan sebagai perilaku yang diamati, bukan langsung dianggap sebagai keberhasilan atau kegagalan.

\section{Lingkungan Penelitian}

Penelitian akan dilakukan pada laptop lokal dengan CPU AMD Ryzen 3 5300U (empat inti dan delapan thread). Sistem memiliki RAM sekitar 9,1 GiB dan grafis Radeon terintegrasi. Konfigurasi penelitian menggunakan CPU, presisi FP32 (bilangan titik mengambang 32-bit), dan ukuran batch satu. Semua metode akan dievaluasi pada perangkat dan pengaturan yang sama.

Versi perangkat lunak, pengaturan thread, revisi model, dan batas token akan dicatat sebelum evaluasi. Pengukuran waktu dilakukan dengan laptop terhubung ke daya, pengaturan daya yang sama, dan aktivitas latar belakang seminimal mungkin. Kesimpulan efisiensi dibatasi pada lingkungan yang diuji.

Rancangan penelitian yang digunakan meliputi pembagian tujuh kelompok, keputusan biner untuk lima kelompok tengah, representasi prompt, arsitektur router, dua tahap pelatihan router dan LoRA, konfigurasi awal adapter, objective beserta kandidat penalti, seed pemilihan subset dan seed utama pelatihan, pemilihan pola static pada data validasi, evaluasi empat varian, perbandingan efisiensi berdasarkan waktu aktual, mode non-thinking, batas token, konfigurasi ROUGE dan BERTScore, tiga kali pemanasan, serta lima kali pengulangan pengukuran.
